\chapter{Conclusiones}
\label{cap:conclusiones}

El presente Trabajo de Fin de Grado ha consistido en diseñar, implementar y validar un sistema de control térmico local que no depende de servicios externos para ninguna de sus funciones críticas, manteniendo los datos del hogar bajo el control del usuario.

A lo largo del proceso de desarrollo, el proyecto ha evolucionado de forma significativa. La fase inicial, apoyada en la infraestructura universitaria para validar el enfoque agéntico, fue necesaria para sentar las bases técnicas y arquitectónicas. La decisión de migrar posteriormente a una solución 100\% local fue el paso lógico para cumplir con el objetivo de eliminar dependencias externas.

\section{Logros Alcanzados}

A fecha de redacción de este capítulo, el proyecto Luxe Core AI cuenta con los siguientes hitos de implementación verificados sobre hardware real:

\begin{itemize}
    \item \textbf{Arquitectura de tres capas escalable:} el orquestador Node.js (OpenClaw npm) coordina los agentes de IA, los canales de usuario y la ejecución de \textit{drivers} Python de dispositivos, cada capa con interfaces bien definidas y gestionables de forma aislada.

    \item \textbf{Control local de dispositivos Tuya (RF-03 verificado):} el enchufe inteligente responde a comandos de activación y desactivación en menos de 2 segundos a través de la red air-gapped, utilizando el protocolo 3.5 de \texttt{tinytuya} sin pasar por los servidores de Tuya. La \textit{local key} se almacena exclusivamente en local. El procedimiento de \textit{Hotspot Spoofing} para el \textit{onboarding} de nuevos dispositivos está documentado y es reproducible.

    \item \textbf{Lectura de sensores ambientales BLE (RF-02 verificado):} los sensores Xiaomi LYWSD03MMC entregan temperatura, humedad y voltaje de batería mediante conexión GATT directa desde el servidor, sin necesidad de flashear firmware alternativo ni de obtener \textit{bind keys} de la nube de Xiaomi. El protocolo de lectura, basado en la característica \texttt{ebe0ccc1} del servicio propietario de Xiaomi, ha sido completamente documentado e implementado en el módulo \texttt{ble\_sensors.py}.

    \item \textbf{Control completo del ventilador FanLamp F8808 (RF-04 verificado):} se ha completado la ingeniería inversa del protocolo propietario Tencent, descubriendo que el ventilador utiliza anuncios BLE (y no conexiones GATT) para recibir comandos. Los anuncios los emite un ESP32 DevKit V1 como radio BLE externa dedicada, controlado desde el servidor mediante \texttt{fanlamp\_control.py} a través del puerto serie USB (115200~baud). Este diseño libera el adaptador Bluetooth del servidor para la lectura simultánea del sensor Xiaomi, elimina la dependencia de \texttt{sudo} y reduce la latencia de $\sim$3~s a $\sim$1,5~s por comando. Se han mapeado y verificado 11 comandos (apagado general, ventilador on/off y 5 velocidades, luz on/off y modo noche). El coste económico de la integración ha sido de 0~\euro{}, reutilizando el ESP32 ya disponible en el proyecto.

    \item \textbf{Orquestador OpenClaw desplegado:} el \textit{gateway} Node.js (OpenClaw npm v2026.5.27) se ejecuta como servicio \texttt{systemd} sobre el puerto 18789, accesible desde toda la subred \texttt{192.168.1.0/24}. Los agentes de IA utilizan Ollama local con los modelos \texttt{qwen2.5-coder:1.5b} (clasificador) y \texttt{mistral:7b-instruct-q4\_K\_M} (razonador).

    \item \textbf{Canales de usuario operativos:} el \textit{gateway} de OpenClaw proporciona integración nativa con Telegram (BotFather API) y un \textit{dashboard} web accesible en \texttt{http://192.168.1.50:18789}, sin necesidad de procesos Python auxiliares.

    \item \textbf{Infraestructura de red aislada validada:} el router TP-Link autónomo con subred \texttt{192.168.1.0/24} y sin salida a Internet garantiza que ningún paquete generado por los dispositivos comerciales abandone el perímetro local. Esta arquitectura ha sido verificada operativamente durante todas las pruebas de integración.

    \item \textbf{Sistema de Zero-Inference masivo (462 patrones):} el Tier~0 del Model Router v2 se expandió de 127 a 462 patrones, cubriendo aproximadamente el 92\% de los comandos diarios con latencia inferior a 1~ms en RAM. Se incluyeron variantes con muletillas, dialectalismos andaluces y formas abreviadas, además de la cobertura completa de iluminación, ventilador, escenas y consultas.

    \item \textbf{Sistema de invalidación de caché:} se implementó un mecanismo de invalidación de la caché LRU del router que permite forzar el reanálisis de comandos cuando se actualizan los patrones o la configuración del sistema, sin necesidad de reiniciar el servicio.

    \item \textbf{Filtro de \texttt{log} de arranque ESP32:} se desarrolló un script que filtra y elimina los mensajes de \texttt{log} de arranque del ESP32 (\textit{boot log}), permitiendo una comunicación más limpia y fiable entre el servidor y el microcontrolador.

    \item \textbf{Disponibilidad del sistema:} durante el período de pruebas continuas, los servicios críticos (model-router, sensor-daemon y openclaw-gateway) permanecieron operativos sin experimentar caídas que requiriesen intervención manual. El sistema se mantuvo en funcionamiento ininterrumpido durante más de 24~horas en el servidor local, confirmando la estabilidad de la arquitectura y la eficacia de los mecanismos de auto-recuperación ante fallos transitorios.
\end{itemize}

\section{Estado Actual del Desarrollo}

\begin{table}[ht]
\centering
\begin{tabular}{|l|c|p{6cm}|}
\hline
\textbf{Módulo} & \textbf{Estado} & \textbf{Observaciones} \\
\hline
Tuya Smart Plug & Completado & Validado en red air-gapped. ON/OFF alterna. \\
BLE Sensors (Xiaomi) & Completado & Lee 25,07~°C / 59~\% / 3055~mV por GATT. \\
FanLamp Pro F8808 & Completado & Controlado vía BLE advertisements con 11 comandos verificados. \\
OpenClaw Orchestrator & Desplegado & \textit{Gateway} Node.js en puerto 18789 (LAN). Canales Telegram + web operativos. \\
Model Router v2 & Desplegado & Tres niveles: Tier~0 (462 patrones, \textasciitilde{}92\% cobertura), Tier~1 (1.5B clasificador), Tier~2 (Mistral 7B razonador). Caché LRU con invalidación. \\
ESP32 HVAC Bridge & Pendiente & Firmware C++20 no implementado aún. \\
\hline
\end{tabular}
\caption{Estado de los módulos del proyecto a la fecha de esta memoria.}
\end{table}

\section{Líneas de Trabajo Futuras}

Las siguientes líneas de trabajo están ordenadas por prioridad y esfuerzo estimado.

\subsection*{Corto plazo (semanas)}
\begin{itemize}
    \item \textbf{Desarrollo del firmware HVAC (ESP32 + LIN):} constituye la prioridad inmediata. El hardware está adquirido y cableado, y el repositorio \texttt{esphome-fujitsu-halcyon} de Omniflux~\cite{omniflux_esphome_fujitsu} proporciona una base sólida que nos ahorrará meses de ingeniería inversa del protocolo LIN. Durante el TFG no se completó por una combinación de alcance y logística doméstica: integrar y validar el firmware requiere poner la máquina fuera de servicio durante días, algo complicado en junio cuando el termómetro en Córdoba supera los 38~$^\circ$C y el aire acondicionado es más necesario que nunca. Queda como primer objetivo tras la entrega de la memoria. \hfill \textbf{(2--3 semanas)}

    \item \textbf{Eliminación del puente ESP32 para el ventilador:} el firmware del ESP32 utilizado como puente BLE presenta reinicios ocasionales que el sistema gestiona mediante reintentos automáticos. Una migración al control BLE directo desde el servidor usando \texttt{bleak}, sin el microcontrolador intermedio, eliminaría este punto único de fallo y reduciría la latencia de los comandos del ventilador. \hfill \textbf{(1--2 semanas)}

    \item \textbf{\textit{Streaming} de respuestas del Tier~2:} actualmente el razonador produce la respuesta completa antes de enviarla al usuario, resultando en 27--60~s de espera. La implementación de \textit{streaming} SSE permitiría mostrar las palabras conforme se generan, mejorando significativamente la experiencia de usuario durante las consultas conversacionales. \hfill \textbf{(1--2 semanas)}

    \item \textbf{Integración del bucle de control con OpenClaw:} conectar los módulos de sensores BLE, actuadores Tuya y control del ventilador FanLamp al bucle principal del orquestador para que el sistema tome decisiones de confort en tiempo real. \hfill \textbf{(1--2 semanas)}

    \item \textbf{Containerización de servicios:} empaquetar los servicios (OpenClaw, Model Router, Sensor Daemon) en contenedores Docker y orquestarlos con \texttt{docker-compose}, facilitando el despliegue en otros entornos y la replicabilidad del proyecto. \hfill \textbf{(1--2 semanas)}
\end{itemize}

\subsection*{Medio plazo (meses)}
\begin{itemize}
    \item \textbf{Expansión del \textit{dashboard}:} extender el \textit{dashboard} nativo de OpenClaw con visualización de analíticas energéticas e histórico de confort ASHRAE 55~\cite{ashrae55_2023}. \hfill \textbf{(1--2 meses)}

    \item \textbf{Expansión sensorial:} integrar un sensor de movimiento por infrarrojos o radar (como el LD2410) para mejorar la detección de presencia sin depender del teléfono móvil. También se plantea incorporar un dispositivo con micrófono y altavoz que permita la comunicación por voz. \hfill \textbf{(1--2 meses)}

    \item \textbf{Base de datos vectorial:} cuando el proyecto escale a decenas de dispositivos y cientos de comandos, incorporar una base de datos vectorial (Qdrant o ChromaDB) para mejorar la recuperación semántica. \hfill \textbf{(2--3 meses)}

    \item \textbf{Memoria conversacional:} se desarrolló un prototipo de memoria de corto plazo basado en \textit{embeddings} con \texttt{bge-m3} que demostró capacidad para detectar referencias implícitas entre mensajes consecutivos. Queda como trabajo futuro su integración estable en el Tier~2, evitando la interferencia con el sistema de patrones del Tier~0 que se observó en el prototipo inicial. \hfill \textbf{(2--3 meses)}
\end{itemize}

\subsection*{Largo plazo (investigación)}
\begin{itemize}
    \item \textbf{Aprendizaje adaptativo:} incorporar modelos de \textit{reinforcement learning} para que el sistema aprenda los hábitos del usuario y anticipe sus preferencias de confort, siguiendo la línea iniciada con el autoaprendizaje del Smart Advisor. Requiere investigación continuada. \hfill \textbf{(varios meses)}

    \item \textbf{Expansión a nuevos protocolos:} integrar dispositivos que usen Zigbee o Z-Wave para ampliar el catálogo de hardware compatible. \hfill \textbf{(varios meses)}

    \item \textbf{Re-evaluación de modelos de lenguaje:} la arquitectura \textit{model-agnostic} del enrutador permite incorporar nuevos modelos modificando únicamente dos líneas en \texttt{config.py}. La familia Qwen3.5, descartada en esta iteración por su modo \textit{thinking} incompatible con CPU, podría volver a evaluarse cuando surjan implementaciones que permitan desactivar este mecanismo. Modelos emergentes como las familias Llama~4 o DeepSeek locales podrían integrarse sin alterar el router. \hfill \textbf{(continuo)}

    \item \textbf{\textit{Fine-tuning} del clasificador:} el Tier~1 emplea \texttt{qwen2.5-coder:1.5b} con un \textit{prompt} de más de 130 ejemplos. Un ajuste fino con datos reales de uso doméstico recogidos durante el despliegue permitiría eliminar el \textit{prompt} y potencialmente mejorar la precisión en español coloquial y dialectalismos andaluces, reduciendo además la latencia al acortar el contexto de entrada. \hfill \textbf{(varios meses)}
\end{itemize}

\section{Reflexión Personal}

Más allá de los logros técnicos, este proyecto ha sido un proceso de aprendizaje continuo. Cada obstáculo (desde un puerto Ethernet mal conectado hasta un protocolo BLE sin documentación) se convirtió en una oportunidad para aprender algo nuevo. He disfrutado especialmente la fase de ingeniería inversa del ventilador FanLamp, donde la combinación de análisis de código, experimentación con hardware y paciencia permitió descifrar un protocolo que el fabricante no había previsto que nadie entendiera.

La decisión de mantener todo el sistema en local, sin depender de la nube, no fue solo técnica: fue también una cuestión de principios. La tecnología del hogar no debería exigir ceder la privacidad a cambio de comodidad. En este proyecto se ha comprobado que es posible tener ambas cosas con hardware de consumo y un enfoque de diseño centrado en el control local.

Los modelos pequeños que utiliza el sistema (1.5B y Mistral 7B) no son superinteligentes. Necesitan el apoyo del Tier~0 para ser útiles en el día a día y, desde luego, no compiten con la fluidez de un Gemini Flash o un GPT-4o en el chat. Pero ponerlos a prueba en un caso real, con hardware limitado y cero dependencia de la nube, me ha convencido de que su potencial es inmenso y que apenas estamos viendo el principio. La evolución de estos modelos en los últimos años ha sido tan rápida que no parece descabellado pensar que, en cuestión de pocos años, un modelo ligero ejecutándose en un procesador doméstico ofrecerá una capacidad de razonamiento cercana a la que hoy solo encontramos en los grandes modelos cloud. Cuando eso ocurra, y creo que ocurrirá, proyectos como este cobrarán aún más sentido: el hardware ya estará ahí, la arquitectura ya estará probada, y la pregunta ya no será «¿se puede hacer en local?», sino «¿por qué harías esto en la nube?».
